\section{Policy language summary}
\label{app:policy}
A pack has an identifier, a version, tiers (decision, description, reversibility), a default tier, tool rules (globs, tier, optional \texttt{when} conditions), argument rules (tools, argument globs, regular expression, \verdict{ask} or \verdict{deny}), clamps, redaction, scope (target arguments, allowed networks and host globs, \texttt{ask}, \texttt{deny} or \texttt{project} on violation), budgets (tools, optional numeric argument, maximum, reserve), typed actions (tools, target, port and rate arguments, limits, schemes, forbidden networks and hosts, wildcard handling, verdict for unreadable values), a predictive section (effects per tool, blast-radius limits, verdict for unknown targets), a flow section (labels by tool, optional argument pattern, destination arguments and internal destinations; rules with a sink and required labels; labels for unclassified tools), a session call limit and the number of denials after which allowed calls need approval. Packs are validated before use; an invalid pack is rejected, never partially applied.

\section{Replay details}
\label{app:replay}
Exposure is found by replacing each injection slot with a unique marker string and executing each user task's ground-truth calls in AgentDojo's runtime; the first call whose returned value contains a marker ends the prefix. Attack goals contribute their ground-truth calls verbatim. The trajectories, with the AgentDojo version and a digest, are part of the artifact, so the evaluation can be rerun without AgentDojo installed. Policies are regenerated from the recorded tool descriptions on every run.
